% operadores_vecindario.tex
% Sección autocontenida que documenta los 9 operadores de vecindario
% del proyecto MetaCARP.  Pensada para incluirse con \input{} desde un
% documento maestro que ya cargue: amsmath, amssymb, algorithm2e
% (ruled/vlined/linesnumbered), hyperref, xcolor, tikz.
%
% NO agregar \documentclass, \begin{document}, etc. aquí.
%
% NOTA SOBRE LA REVISIÓN (2026): los pseudocódigos de esta versión son
% AUTOCONTENIDOS y fieles a metacarp/vecindarios.py.  A diferencia de la
% versión anterior, cada algoritmo muestra explícitamente los cinco pasos
% que el código real ejecuta alrededor del movimiento puro y que los
% pseudocódigos ingenuos omitían:
%   (1) Normalización: se quita el token de depósito "D" antes de operar.
%   (2) Filtrado de rutas candidatas según el requisito del operador.
%   (3) Bucle de reintentos (hasta 500): ante un sorteo inaplicable se hace
%       "continue" en vez de devolver null.
%   (4) Copia profunda de la solución antes de mutarla (Python pasa listas
%       por referencia).
%   (5) Desnormalización: se restaura "D" al inicio y fin de cada ruta.

% ============================================================
\section{Operadores de Vecindario}
\label{sec:neighborhood-operators}
% ============================================================

Un \emph{operador de vecindario} define una transformación pequeña y
reversible que convierte una solución factible del CARP $s$ en un
\emph{vecino} $s'$.
En cada iteración, todas las metaheurísticas de MetaCARP (Búsqueda Tabú,
Recocido Simulado, Colonia de Abejas, Cuckoo Search) llaman a una única
rutina maestra que elige un operador al azar ---con pesos no uniformes
opcionales--- y lo aplica a la solución actual.  Nueve operadores están
registrados en \texttt{OPERADORES\_POPULARES}; se dividen en dos familias:

\begin{itemize}
    \item \textbf{Intra-ruta} (3 operadores): actúan sobre una sola ruta
          $r_a$; el conjunto de tareas asignado a cada ruta no cambia.
    \item \textbf{Inter-ruta} (6 operadores): transfieren o intercambian
          tareas entre dos rutas distintas $r_a$ y $r_b$; pueden reequilibrar
          la carga entre rutas y se prefieren cuando se detecta una violación
          de capacidad.
\end{itemize}

% ------------------------------------------------------------
\subsection{Notación y andamiaje común}
\label{subsec:notation-ops}
% ------------------------------------------------------------

A lo largo de esta sección rigen las siguientes convenciones.

\begin{itemize}
    \item $r[i]$: tarea en la posición 0-indexada $i$ de la ruta $r$.
    \item $r[i{:}j]$: subsecuencia contigua desde la posición $i$ hasta la
          posición $j$ inclusive (ambos extremos incluidos).
    \item $r_a$, $r_b$: dos rutas distintas ($a \neq b$), salvo en el caso
          de respaldo intra-ruta de \texttt{or\_opt\_2} / \texttt{or\_opt\_3}.
    \item $\mathit{head}(r, c) = r[0{:}c]$,\enspace
          $\mathit{tail}(r, c) = r[c{+}1{:}]$.
    \item $\|$ denota concatenación de listas.
    \item Las etiquetas de tarea se escriben $\mathtt{TR1}, \mathtt{TR2},
          \ldots$; el token de depósito \texttt{D} se elimina antes de
          aplicar cualquier operador y se restaura después.
\end{itemize}

\paragraph{Pasos compartidos por todos los operadores.}
Los pseudocódigos de esta sección son \emph{autocontenidos}: cada uno refleja
exactamente lo que ejecuta \texttt{generar\_vecino} en
\texttt{metacarp/vecindarios.py}.  Cinco pasos rodean al movimiento puro y
aparecen en \textbf{todos} los algoritmos:

\begin{enumerate}
    \item \textsc{NormalizarQuitarDepósito}$(s)$: elimina el token \texttt{D}
          del inicio y fin de cada ruta, dejando solo etiquetas de tarea
          ($\hat{s}$).  Sin este paso, un \texttt{pop} podría extraer un
          \texttt{D} en lugar de una tarea real.
    \item Filtrado de \emph{rutas candidatas} según el requisito mínimo del
          operador (p.\,ej.\ $|r| \geq 2$ para \texttt{relocate\_intra}).
    \item Bucle \textsc{Repeat} de hasta $T_{\max}=500$ intentos: ante un
          sorteo inaplicable se ejecuta \textbf{continue} y se vuelve a
          intentar; si se supera $T_{\max}$ se lanza \texttt{RuntimeError}.
          (Donde un pseudocódigo ingenuo escribía \texttt{Return null}, el
          código real hace \textbf{continue}.)
    \item \textsc{CopiaProfunda}$(\hat{s})$: Python pasa las listas por
          referencia; sin una copia independiente el movimiento corrompería
          la solución original.
    \item \textsc{RestaurarDepósito}$(\hat{s}')$: vuelve a añadir \texttt{D}
          al inicio y fin de cada ruta antes de devolver el vecino, junto
          con un registro \textsc{MovimientoVecindario} que describe el
          movimiento aplicado.
\end{enumerate}

% ============================================================
\subsection{Operadores Intra-ruta}
\label{subsec:intra-operators}
% ============================================================

% ---- relocate_intra ----
\subsubsection{\texttt{relocate\_intra}}
\label{sssec:relocate-intra}

\textbf{Movimiento.}
Quita la tarea en la posición $i$ de la ruta $r_a$ y la reinserta en la
posición $j \neq i$ de la misma ruta.

\textbf{Requisito.} $|r_a| \geq 2$.

\begin{algorithm}[H]
\DontPrintSemicolon
\caption{\texttt{relocate\_intra} --- Mover una tarea a otra posición dentro de la misma ruta}
\label{alg:relocate-intra}

\KwIn{Solución con depósito $s$;\; generador aleatorio $\mathit{rng}$;\; máximo de intentos $T_{\max}=500$}
\KwOut{Vecino con depósito $s'$;\; registro de movimiento $m$}

\BlankLine
$\hat{s} \leftarrow \textsc{NormalizarQuitarDepósito}(s)$
\tcp{elimina el token "D" de inicio y fin de cada ruta}
$t \leftarrow 0$\;
\Repeat{movimiento válido encontrado}{
    $t \leftarrow t + 1$\;
    \lIf{$t > T_{\max}$}{\textbf{lanzar} \texttt{RuntimeError}}
    $\mathcal{C} \leftarrow \{k : |\hat{r}_k| \geq 2\}$
    \tcp{rutas candidatas: al menos 2 tareas}
    \lIf{$\mathcal{C} = \emptyset$}{\textbf{continue}}
    $a \leftarrow \mathit{rng}.\textsc{choice}(\mathcal{C})$;\quad $n \leftarrow |\hat{r}_a|$\;
    $i \leftarrow \mathit{rng}.\textsc{randrange}(n)$
    \tcp{posición de origen, uniforme en $\{0,\ldots,n-1\}$}
    $j \leftarrow \mathit{rng}.\textsc{randrange}(n)$
    \tcp{posición de destino, uniforme en $\{0,\ldots,n-1\}$}
    \lIf{$i = j$}{\textbf{continue} \tcp*{no genera vecino distinto}}
    $\hat{s}' \leftarrow \textsc{CopiaProfunda}(\hat{s})$
    \tcp{Python pasa listas por referencia; sin copia se corrompería $\hat{s}$}
    $e \leftarrow \hat{r}'_a.\textsc{pop}(i)$\;
    $\hat{r}'_a.\textsc{insert}(j,\, e)$\;
    $m \leftarrow \textsc{MovimientoVecindario}(\texttt{relocate\_intra},\; \mathit{ruta\_a}=a,\; i=i,\; j=j)$\;
    $s' \leftarrow \textsc{RestaurarDepósito}(\hat{s}')$
    \tcp{añade "D" al inicio y fin de cada ruta antes de devolver}
    \Return $(s',\, m)$\;
}
\end{algorithm}

\noindent\textbf{Ejemplo visual} ($i=0$, $j=2$):
\[
    r_a = [\,\mathtt{TR1},\; \mathtt{TR2},\; \mathtt{TR3},\; \mathtt{TR4}\,]
    \;\xrightarrow{\;\text{pop}(0),\;\text{insert}(2)\;}\;
    r'_a = [\,\mathtt{TR2},\; \mathtt{TR3},\; \mathtt{TR1},\; \mathtt{TR4}\,]
\]

% ---- swap_intra ----
\subsubsection{\texttt{swap\_intra}}
\label{sssec:swap-intra}

\textbf{Movimiento.}
Intercambia las tareas en las posiciones $i$ y $j$ ($i \neq j$) dentro de
la misma ruta $r_a$.

\textbf{Requisito.} $|r_a| \geq 2$.

\begin{algorithm}[H]
\DontPrintSemicolon
\caption{\texttt{swap\_intra} --- Intercambiar dos tareas dentro de la misma ruta}
\label{alg:swap-intra}

\KwIn{Solución con depósito $s$;\; generador aleatorio $\mathit{rng}$;\; máximo de intentos $T_{\max}=500$}
\KwOut{Vecino con depósito $s'$;\; registro de movimiento $m$}

\BlankLine
$\hat{s} \leftarrow \textsc{NormalizarQuitarDepósito}(s)$
\tcp{elimina el token "D" de inicio y fin de cada ruta}
$t \leftarrow 0$\;
\Repeat{movimiento válido encontrado}{
    $t \leftarrow t + 1$\;
    \lIf{$t > T_{\max}$}{\textbf{lanzar} \texttt{RuntimeError}}
    $\mathcal{C} \leftarrow \{k : |\hat{r}_k| \geq 2\}$
    \tcp{rutas candidatas: al menos 2 tareas}
    \lIf{$\mathcal{C} = \emptyset$}{\textbf{continue}}
    $a \leftarrow \mathit{rng}.\textsc{choice}(\mathcal{C})$;\quad $n \leftarrow |\hat{r}_a|$\;
    $(i,\, j) \leftarrow \mathit{rng}.\textsc{sample}(\{0,\ldots,n-1\},\, 2)$
    \tcp{dos índices distintos sin reemplazo (garantiza $i \neq j$ por construcción)}
    $\hat{s}' \leftarrow \textsc{CopiaProfunda}(\hat{s})$
    \tcp{Python pasa listas por referencia; sin copia se corrompería $\hat{s}$}
    $\hat{r}'_a[i],\, \hat{r}'_a[j] \leftarrow \hat{r}'_a[j],\, \hat{r}'_a[i]$
    \tcp{intercambio atómico: lado derecho se evalúa completo antes de asignar}
    $m \leftarrow \textsc{MovimientoVecindario}(\texttt{swap\_intra},\; \mathit{ruta\_a}=a,\; i=i,\; j=j)$\;
    $s' \leftarrow \textsc{RestaurarDepósito}(\hat{s}')$
    \tcp{añade "D" al inicio y fin de cada ruta antes de devolver}
    \Return $(s',\, m)$\;
}
\end{algorithm}

\noindent\textbf{Ejemplo visual} ($i=0$, $j=3$):
\[
    r_a = [\,\mathtt{TR1},\; \mathtt{TR2},\; \mathtt{TR3},\; \mathtt{TR4}\,]
    \;\xrightarrow{\;\text{swap}(0,3)\;}\;
    r'_a = [\,\mathtt{TR4},\; \mathtt{TR2},\; \mathtt{TR3},\; \mathtt{TR1}\,]
\]

% ---- 2opt_intra ----
\subsubsection{\texttt{2opt\_intra}}
\label{sssec:2opt-intra}

\textbf{Movimiento.}
Invierte la subsecuencia contigua $r_a[i{:}j]$ (ambos extremos inclusive),
con $i < j$.  Las tareas fuera del segmento no se ven afectadas.

\textbf{Requisito.} $|r_a| \geq 3$ (garantiza un segmento de al menos dos
tareas cuando $j > i$).

\begin{algorithm}[H]
\DontPrintSemicolon
\caption{\texttt{2opt\_intra} --- Invertir un segmento de tareas dentro de la misma ruta}
\label{alg:2opt-intra}

\KwIn{Solución con depósito $s$;\; generador aleatorio $\mathit{rng}$;\; máximo de intentos $T_{\max}=500$}
\KwOut{Vecino con depósito $s'$;\; registro de movimiento $m$}

\BlankLine
$\hat{s} \leftarrow \textsc{NormalizarQuitarDepósito}(s)$
\tcp{elimina el token "D" de inicio y fin de cada ruta}
$t \leftarrow 0$\;
\Repeat{movimiento válido encontrado}{
    $t \leftarrow t + 1$\;
    \lIf{$t > T_{\max}$}{\textbf{lanzar} \texttt{RuntimeError}}
    $\mathcal{C} \leftarrow \{k : |\hat{r}_k| \geq 3\}$
    \tcp{rutas candidatas: mínimo 3 tareas para garantizar segmento de longitud $\geq 2$}
    \lIf{$\mathcal{C} = \emptyset$}{\textbf{continue}}
    $a \leftarrow \mathit{rng}.\textsc{choice}(\mathcal{C})$;\quad $n \leftarrow |\hat{r}_a|$\;
    $i \leftarrow \mathit{rng}.\textsc{randrange}(0,\, n-1)$
    \tcp{inicio del segmento, uniforme en $\{0,\ldots,n-2\}$}
    $j \leftarrow \mathit{rng}.\textsc{randrange}(i+1,\, n)$
    \tcp{fin del segmento, uniforme en $\{i+1,\ldots,n-1\}$; garantiza $j > i$}
    $\hat{s}' \leftarrow \textsc{CopiaProfunda}(\hat{s})$
    \tcp{Python pasa listas por referencia; sin copia se corrompería $\hat{s}$}
    $\hat{r}'_a[i\,{:}\,j{+}1] \leftarrow \textsc{Revertir}(\hat{r}'_a[i\,{:}\,j{+}1])$
    \tcp{asignación de slice: reemplaza el segmento en su lugar}
    $m \leftarrow \textsc{MovimientoVecindario}(\texttt{2opt\_intra},\; \mathit{ruta\_a}=a,\; i=i,\; j=j)$\;
    $s' \leftarrow \textsc{RestaurarDepósito}(\hat{s}')$
    \tcp{añade "D" al inicio y fin de cada ruta antes de devolver}
    \Return $(s',\, m)$\;
}
\end{algorithm}

\noindent\textbf{Ejemplo visual} ($i=1$, $j=3$):
\[
    r_a =
    \bigl[\,\mathtt{TR1},\;
    \underbrace{\mathtt{TR2},\; \mathtt{TR3},\; \mathtt{TR4}}_{i=1,\;j=3},\;
    \mathtt{TR5}\,\bigr]
    \;\longrightarrow\;
    r'_a = \bigl[\,\mathtt{TR1},\; \mathtt{TR4},\; \mathtt{TR3},\; \mathtt{TR2},\; \mathtt{TR5}\,\bigr]
\]

% ============================================================
\subsection{Operadores Inter-ruta}
\label{subsec:inter-operators}
% ============================================================

% ---- relocate_inter ----
\subsubsection{\texttt{relocate\_inter}}
\label{sssec:relocate-inter}

\textbf{Movimiento.}
Quita la tarea en la posición $i$ de la ruta $r_a$ y la inserta en la
posición $j$ de la ruta $r_b$, con $a \neq b$.

\textbf{Requisito.} $|r_a| \geq 1$; $r_b$ puede estar vacía (la tarea
simplemente se agrega).

\begin{algorithm}[H]
\DontPrintSemicolon
\caption{\texttt{relocate\_inter} --- Mover una tarea de una ruta a otra ruta distinta}
\label{alg:relocate-inter}

\KwIn{Solución con depósito $s$;\; generador aleatorio $\mathit{rng}$;\; máximo de intentos $T_{\max}=500$}
\KwOut{Vecino con depósito $s'$;\; registro de movimiento $m$}

\BlankLine
$\hat{s} \leftarrow \textsc{NormalizarQuitarDepósito}(s)$
\tcp{elimina el token "D" de inicio y fin de cada ruta}
$t \leftarrow 0$\;
\Repeat{movimiento válido encontrado}{
    $t \leftarrow t + 1$\;
    \lIf{$t > T_{\max}$}{\textbf{lanzar} \texttt{RuntimeError}}
    $\mathcal{A} \leftarrow \{k : |\hat{r}_k| \geq 1\}$
    \tcp{rutas con al menos 1 tarea (candidatas como origen)}
    \lIf{$|\mathcal{A}| < 1$ \textbf{o} $|\hat{s}| < 2$}{\textbf{continue}}
    $\mathit{ra} \leftarrow \mathit{rng}.\textsc{choice}(\mathcal{A})$
    \tcp{ruta origen: uniforme entre rutas no vacías}
    $\mathit{rb} \leftarrow \mathit{rng}.\textsc{randrange}(|\hat{s}|)$
    \tcp{ruta destino: uniforme entre \textbf{todas} las rutas (puede estar vacía)}
    \lIf{$\mathit{ra} = \mathit{rb}$}{\textbf{continue} \tcp*{debe ser inter-ruta}}
    $i \leftarrow \mathit{rng}.\textsc{randrange}(|\hat{r}_{\mathit{ra}}|)$
    \tcp{posición en la ruta origen}
    $j \leftarrow \mathit{rng}.\textsc{randrange}(|\hat{r}_{\mathit{rb}}| + 1)$
    \tcp{posición en la ruta destino; $+1$ permite insertar al final}
    $\hat{s}' \leftarrow \textsc{CopiaProfunda}(\hat{s})$
    \tcp{Python pasa listas por referencia; sin copia se corrompería $\hat{s}$}
    $e \leftarrow \hat{r}'_{\mathit{ra}}.\textsc{pop}(i)$\;
    $\hat{r}'_{\mathit{rb}}.\textsc{insert}(j,\, e)$\;
    $m \leftarrow \textsc{MovimientoVecindario}(\texttt{relocate\_inter},\; \mathit{ruta\_a}=\mathit{ra},\; \mathit{ruta\_b}=\mathit{rb},\; i=i,\; j=j)$\;
    $s' \leftarrow \textsc{RestaurarDepósito}(\hat{s}')$
    \tcp{añade "D" al inicio y fin de cada ruta antes de devolver}
    \Return $(s',\, m)$\;
}
\end{algorithm}

\noindent\textbf{Ejemplo visual} ($i=2$, $j=1$):
\[
    \begin{array}{ll}
        r_a = [\mathtt{TR1},\;\mathtt{TR2},\;\mathtt{TR3}]
        & r_b = [\mathtt{TR4},\;\mathtt{TR5}] \\[4pt]
        \multicolumn{2}{c}{\xrightarrow{\;\text{mover}\;\mathtt{TR3}\;\text{de}\;r_a[2]\;\text{a}\;r_b[1]\;}} \\[4pt]
        r'_a = [\mathtt{TR1},\;\mathtt{TR2}]
        & r'_b = [\mathtt{TR4},\;\mathtt{TR3},\;\mathtt{TR5}]
    \end{array}
\]

% ---- swap_inter ----
\subsubsection{\texttt{swap\_inter}}
\label{sssec:swap-inter}

\textbf{Movimiento.}
Intercambia la tarea en la posición $i$ de la ruta $r_a$ con la tarea en
la posición $j$ de la ruta $r_b$ ($a \neq b$).  Ambas rutas conservan el
mismo número de tareas.

\textbf{Requisito.} $|r_a| \geq 1$ y $|r_b| \geq 1$.

\begin{algorithm}[H]
\DontPrintSemicolon
\caption{\texttt{swap\_inter} --- Intercambiar una tarea entre dos rutas distintas}
\label{alg:swap-inter}

\KwIn{Solución con depósito $s$;\; generador aleatorio $\mathit{rng}$;\; máximo de intentos $T_{\max}=500$}
\KwOut{Vecino con depósito $s'$;\; registro de movimiento $m$}

\BlankLine
$\hat{s} \leftarrow \textsc{NormalizarQuitarDepósito}(s)$
\tcp{elimina el token "D" de inicio y fin de cada ruta}
$t \leftarrow 0$\;
\Repeat{movimiento válido encontrado}{
    $t \leftarrow t + 1$\;
    \lIf{$t > T_{\max}$}{\textbf{lanzar} \texttt{RuntimeError}}
    \lIf{$|\hat{s}| < 2$}{\textbf{continue}}
    $\mathcal{N} \leftarrow \{k : |\hat{r}_k| \geq 1\}$
    \tcp{rutas no vacías; ambas rutas participantes deben tener al menos 1 tarea}
    \lIf{$|\mathcal{N}| < 2$}{\textbf{continue}}
    $(\mathit{ra},\, \mathit{rb}) \leftarrow \mathit{rng}.\textsc{sample}(\mathcal{N},\, 2)$
    \tcp{dos rutas distintas no vacías, sin reemplazo; garantiza $\mathit{ra} \neq \mathit{rb}$}
    $i \leftarrow \mathit{rng}.\textsc{randrange}(|\hat{r}_{\mathit{ra}}|)$
    \tcp{posición en la ruta $\mathit{ra}$}
    $j \leftarrow \mathit{rng}.\textsc{randrange}(|\hat{r}_{\mathit{rb}}|)$
    \tcp{posición en la ruta $\mathit{rb}$}
    $\hat{s}' \leftarrow \textsc{CopiaProfunda}(\hat{s})$
    \tcp{Python pasa listas por referencia; sin copia se corrompería $\hat{s}$}
    $\hat{r}'_{\mathit{ra}}[i],\, \hat{r}'_{\mathit{rb}}[j] \leftarrow \hat{r}'_{\mathit{rb}}[j],\, \hat{r}'_{\mathit{ra}}[i]$
    \tcp{intercambio atómico: lado derecho se evalúa completo antes de asignar}
    $m \leftarrow \textsc{MovimientoVecindario}(\texttt{swap\_inter},\; \mathit{ruta\_a}=\mathit{ra},\; \mathit{ruta\_b}=\mathit{rb},\; i=i,\; j=j)$\;
    $s' \leftarrow \textsc{RestaurarDepósito}(\hat{s}')$
    \tcp{añade "D" al inicio y fin de cada ruta antes de devolver}
    \Return $(s',\, m)$\;
}
\end{algorithm}

\noindent\textbf{Ejemplo visual} ($i=0$, $j=1$):
\[
    \begin{array}{ll}
        r_a = [\mathtt{TR1},\;\mathtt{TR2}]
        & r_b = [\mathtt{TR3},\;\mathtt{TR4}] \\[4pt]
        \multicolumn{2}{c}{\xrightarrow{\;\text{swap}\;\mathtt{TR1}\leftrightarrow\mathtt{TR4}\;}} \\[4pt]
        r'_a = [\mathtt{TR4},\;\mathtt{TR2}]
        & r'_b = [\mathtt{TR3},\;\mathtt{TR1}]
    \end{array}
\]

% ---- 2opt_star ----
\subsubsection{\texttt{2opt\_star}}
\label{sssec:2opt-star}

\textbf{Movimiento.}
Divide la ruta $r_a$ en el punto de corte $c_a$ en cabeza $r_a[0{:}c_a]$ y
cola $r_a[c_a{+}1{:}]$; divide $r_b$ en $c_b$ de igual modo.  Reensambla
intercambiando las colas:
\[
    r'_a = \mathit{head}_a \;\|\; \mathit{tail}_b,
    \qquad
    r'_b = \mathit{head}_b \;\|\; \mathit{tail}_a.
\]
Es el análogo inter-ruta del 2-opt clásico; puede corregir rutas que se
cruzan geográficamente.

\textbf{Requisito.} $|r_a| \geq 1$ y $|r_b| \geq 1$.

\begin{algorithm}[H]
\DontPrintSemicolon
\caption{\texttt{2opt\_star} --- Intercambiar las colas de dos rutas (2-opt*)}
\label{alg:2opt-star}

\KwIn{Solución con depósito $s$;\; generador aleatorio $\mathit{rng}$;\; máximo de intentos $T_{\max}=500$}
\KwOut{Vecino con depósito $s'$;\; registro de movimiento $m$}

\BlankLine
$\hat{s} \leftarrow \textsc{NormalizarQuitarDepósito}(s)$
\tcp{elimina el token "D" de inicio y fin de cada ruta}
$t \leftarrow 0$\;
\Repeat{movimiento válido encontrado}{
    $t \leftarrow t + 1$\;
    \lIf{$t > T_{\max}$}{\textbf{lanzar} \texttt{RuntimeError}}
    \lIf{$|\hat{s}| < 2$}{\textbf{continue}}
    $\mathcal{N} \leftarrow \{k : |\hat{r}_k| \geq 1\}$
    \tcp{rutas no vacías; ambas participantes necesitan al menos 1 tarea}
    \lIf{$|\mathcal{N}| < 2$}{\textbf{continue}}
    $(\mathit{ra},\, \mathit{rb}) \leftarrow \mathit{rng}.\textsc{sample}(\mathcal{N},\, 2)$
    \tcp{dos rutas distintas no vacías, sin reemplazo}
    $\mathit{cut\_a} \leftarrow \mathit{rng}.\textsc{randrange}(|\hat{r}_{\mathit{ra}}|)$
    \tcp{punto de corte en la ruta A, uniforme en $\{0,\ldots,|\hat{r}_{\mathit{ra}}|-1\}$}
    $\mathit{cut\_b} \leftarrow \mathit{rng}.\textsc{randrange}(|\hat{r}_{\mathit{rb}}|)$
    \tcp{punto de corte en la ruta B, uniforme en $\{0,\ldots,|\hat{r}_{\mathit{rb}}|-1\}$}
    $\hat{s}' \leftarrow \textsc{CopiaProfunda}(\hat{s})$
    \tcp{Python pasa listas por referencia; sin copia se corrompería $\hat{s}$}
    $\mathit{head\_A} \leftarrow \hat{r}'_{\mathit{ra}}[:\mathit{cut\_a}{+}1]$;\quad
    $\mathit{tail\_A} \leftarrow \hat{r}'_{\mathit{ra}}[\mathit{cut\_a}{+}1:]$\;
    $\mathit{head\_B} \leftarrow \hat{r}'_{\mathit{rb}}[:\mathit{cut\_b}{+}1]$;\quad
    $\mathit{tail\_B} \leftarrow \hat{r}'_{\mathit{rb}}[\mathit{cut\_b}{+}1:]$\;
    $\hat{r}'_{\mathit{ra}} \leftarrow \mathit{head\_A} \mathbin{\|} \mathit{tail\_B}$
    \tcp{ruta A adopta la cola de B}
    $\hat{r}'_{\mathit{rb}} \leftarrow \mathit{head\_B} \mathbin{\|} \mathit{tail\_A}$
    \tcp{ruta B adopta la cola de A}
    $m \leftarrow \textsc{MovimientoVecindario}(\texttt{2opt\_star},\; \mathit{ruta\_a}=\mathit{ra},\; \mathit{ruta\_b}=\mathit{rb},\; i=\mathit{cut\_a},\; j=\mathit{cut\_b})$\;
    $s' \leftarrow \textsc{RestaurarDepósito}(\hat{s}')$
    \tcp{añade "D" al inicio y fin de cada ruta antes de devolver}
    \Return $(s',\, m)$\;
}
\end{algorithm}

\noindent\textbf{Ejemplo visual} ($c_a=1$, $c_b=1$):
\[
    \begin{array}{ll}
        r_a = \bigl[\,\underbrace{\mathtt{TR1},\;\mathtt{TR2}}_{\mathit{head}_a},\;
                     \overbrace{\mathtt{TR3},\;\mathtt{TR4}}^{\mathit{tail}_a}\,\bigr]
        &
        r_b = \bigl[\,\underbrace{\mathtt{TR5},\;\mathtt{TR6}}_{\mathit{head}_b},\;
                     \overbrace{\mathtt{TR7}}^{\mathit{tail}_b}\,\bigr] \\[10pt]
        \multicolumn{2}{c}{\Downarrow} \\[4pt]
        r'_a = [\mathtt{TR1},\;\mathtt{TR2},\;\mathtt{TR7}]
        &
        r'_b = [\mathtt{TR5},\;\mathtt{TR6},\;\mathtt{TR3},\;\mathtt{TR4}]
    \end{array}
\]

% ---- cross_exchange ----
\subsubsection{\texttt{cross\_exchange}}
\label{sssec:cross-exchange}

\textbf{Movimiento.}
Extrae un segmento contiguo $\mathit{seg}_a = r_a[i{:}j]$ de $r_a$ y un
segmento contiguo $\mathit{seg}_b = r_b[k{:}l]$ de $r_b$ (todos los índices
inclusive, $i \leq j$, $k \leq l$), y luego intercambia los segmentos:
\[
    r'_a = r_a[0{:}i] \;\|\; \mathit{seg}_b \;\|\; r_a[j{+}1{:}],
    \qquad
    r'_b = r_b[0{:}k] \;\|\; \mathit{seg}_a \;\|\; r_b[l{+}1{:}].
\]
Los segmentos pueden tener distinta longitud, lo que hace este movimiento
más disruptivo que \texttt{swap\_inter}.

\textbf{Requisito.} $|r_a| \geq 2$ y $|r_b| \geq 2$ (cada segmento debe
contener al menos dos tareas para que puedan cumplirse $j > i$ y $l > k$).

\begin{algorithm}[H]
\DontPrintSemicolon
\caption{\texttt{cross\_exchange} --- Intercambiar segmentos contiguos entre dos rutas}
\label{alg:cross-exchange}

\KwIn{Solución con depósito $s$;\; generador aleatorio $\mathit{rng}$;\; máximo de intentos $T_{\max}=500$}
\KwOut{Vecino con depósito $s'$;\; registro de movimiento $m$}

\BlankLine
$\hat{s} \leftarrow \textsc{NormalizarQuitarDepósito}(s)$
\tcp{elimina el token "D" de inicio y fin de cada ruta}
$t \leftarrow 0$\;
\Repeat{movimiento válido encontrado}{
    $t \leftarrow t + 1$\;
    \lIf{$t > T_{\max}$}{\textbf{lanzar} \texttt{RuntimeError}}
    \lIf{$|\hat{s}| < 2$}{\textbf{continue}}
    $\mathcal{N} \leftarrow \{k : |\hat{r}_k| \geq 2\}$
    \tcp{rutas con al menos 2 tareas en ambas participantes (requisito del segmento)}
    \lIf{$|\mathcal{N}| < 2$}{\textbf{continue}}
    $(\mathit{ra},\, \mathit{rb}) \leftarrow \mathit{rng}.\textsc{sample}(\mathcal{N},\, 2)$
    \tcp{dos rutas distintas con $\geq 2$ tareas, sin reemplazo}
    $n_a \leftarrow |\hat{r}_{\mathit{ra}}|$;\quad $n_b \leftarrow |\hat{r}_{\mathit{rb}}|$\;
    $i \leftarrow \mathit{rng}.\textsc{randrange}(0,\, n_a - 1)$
    \tcp{inicio del segmento en A, uniforme en $\{0,\ldots,n_a-2\}$}
    $j \leftarrow \mathit{rng}.\textsc{randrange}(i+1,\, n_a)$
    \tcp{fin del segmento en A, uniforme en $\{i+1,\ldots,n_a-1\}$; garantiza $j > i$}
    $k \leftarrow \mathit{rng}.\textsc{randrange}(0,\, n_b - 1)$
    \tcp{inicio del segmento en B, uniforme en $\{0,\ldots,n_b-2\}$}
    $l \leftarrow \mathit{rng}.\textsc{randrange}(k+1,\, n_b)$
    \tcp{fin del segmento en B, uniforme en $\{k+1,\ldots,n_b-1\}$; garantiza $l > k$}
    $\hat{s}' \leftarrow \textsc{CopiaProfunda}(\hat{s})$
    \tcp{Python pasa listas por referencia; sin copia se corrompería $\hat{s}$}
    $\mathit{seg\_A} \leftarrow \hat{r}'_{\mathit{ra}}[i\,{:}\,j{+}1]$;\quad
    $\mathit{seg\_B} \leftarrow \hat{r}'_{\mathit{rb}}[k\,{:}\,l{+}1]$\;
    $\hat{r}'_{\mathit{ra}} \leftarrow \hat{r}'_{\mathit{ra}}[{:}i] \mathbin{\|} \mathit{seg\_B} \mathbin{\|} \hat{r}'_{\mathit{ra}}[j{+}1{:}]$\;
    $\hat{r}'_{\mathit{rb}} \leftarrow \hat{r}'_{\mathit{rb}}[{:}k] \mathbin{\|} \mathit{seg\_A} \mathbin{\|} \hat{r}'_{\mathit{rb}}[l{+}1{:}]$\;
    $m \leftarrow \textsc{MovimientoVecindario}(\texttt{cross\_exchange},\; \mathit{ruta\_a}=\mathit{ra},\; \mathit{ruta\_b}=\mathit{rb},\; i=i,\; j=j,\; k=k,\; l=l)$\;
    $s' \leftarrow \textsc{RestaurarDepósito}(\hat{s}')$
    \tcp{añade "D" al inicio y fin de cada ruta antes de devolver}
    \Return $(s',\, m)$\;
}
\end{algorithm}

\noindent\textbf{Ejemplo visual} ($i=1$, $j=2$, $k=0$, $l=1$):
\[
    \begin{array}{ll}
        r_a = \bigl[\mathtt{TR1},\;
                    \overbrace{\mathtt{TR2},\;\mathtt{TR3}}^{\mathit{seg}_a},\;
                    \mathtt{TR4}\bigr]
        &
        r_b = \bigl[\overbrace{\mathtt{TR5},\;\mathtt{TR6}}^{\mathit{seg}_b},\;
                    \mathtt{TR7}\bigr] \\[8pt]
        \multicolumn{2}{c}{\Downarrow} \\[4pt]
        r'_a = [\mathtt{TR1},\;\mathtt{TR5},\;\mathtt{TR6},\;\mathtt{TR4}]
        &
        r'_b = [\mathtt{TR2},\;\mathtt{TR3},\;\mathtt{TR7}]
    \end{array}
\]

% ---- or_opt_2 ----
\subsubsection{\texttt{or\_opt\_2}}
\label{sssec:or-opt-2}

\textbf{Movimiento.}
Extrae el bloque contiguo de exactamente \textbf{2} tareas
$\mathit{block} = r_a[i{:}i{+}2]$ de la ruta $r_a$ y lo inserta en la
posición $j$ de la ruta $r_b$ ($a \neq b$) \emph{en el mismo orden}
(nunca invertido).  Esto distingue a \texttt{or\_opt\_2} de
\texttt{cross\_exchange}, que puede intercambiar segmentos de distinto
tamaño y no preserva el orden de inserción.

\textbf{Requisito.} $|r_a| \geq 2$.  Si la solución tiene una sola ruta,
$r_b$ puede igualar a $r_a$ (respaldo intra): el bloque se quita y se
reinserta en otra posición de la misma ruta.

\begin{algorithm}[H]
\DontPrintSemicolon
\caption{\texttt{or\_opt\_2} --- Mover un bloque de 2 tareas consecutivas a otra ruta}
\label{alg:or-opt-2}

\KwIn{Solución con depósito $s$;\; generador aleatorio $\mathit{rng}$;\; máximo de intentos $T_{\max}=500$}
\KwOut{Vecino con depósito $s'$;\; registro de movimiento $m$}

\BlankLine
$\hat{s} \leftarrow \textsc{NormalizarQuitarDepósito}(s)$
\tcp{elimina el token "D" de inicio y fin de cada ruta}
$t \leftarrow 0$\;
\Repeat{movimiento válido encontrado}{
    $t \leftarrow t + 1$\;
    \lIf{$t > T_{\max}$}{\textbf{lanzar} \texttt{RuntimeError}}
    $\mathcal{C} \leftarrow \{k : |\hat{r}_k| \geq 2\}$
    \tcp{rutas candidatas como origen: al menos 2 tareas (tamaño del bloque)}
    \lIf{$\mathcal{C} = \emptyset$}{\textbf{continue}}
    $\mathit{ra} \leftarrow \mathit{rng}.\textsc{choice}(\mathcal{C})$\;
    $\mathcal{D} \leftarrow \{x : x \neq \mathit{ra}\}$
    \tcp{rutas destino distintas a la origen; si solo hay una, $\mathcal{D} = \{\mathit{ra}\}$}
    \lIf{$\mathcal{D} = \emptyset$}{$\mathcal{D} \leftarrow \{\mathit{ra}\}$}
    $\mathit{rb} \leftarrow \mathit{rng}.\textsc{choice}(\mathcal{D})$\;
    $n_a \leftarrow |\hat{r}_{\mathit{ra}}|$\;
    $i \leftarrow \mathit{rng}.\textsc{randrange}(0,\, n_a - 2 + 1)$
    \tcp{inicio del bloque; garantiza que caben 2 tareas desde $i$}
    \eIf{$\mathit{rb} = \mathit{ra}$}{
        $\tau \leftarrow n_a - 2$
        \tcp{posiciones disponibles tras extraer el bloque}
        \lIf{$\tau \leq 0$}{\textbf{continue}}
        $j \leftarrow \mathit{rng}.\textsc{randrange}(0,\, \tau + 1)$\;
        \lIf{$j = i$}{\textbf{continue} \tcp*{no genera vecino distinto}}
    }{
        $j \leftarrow \mathit{rng}.\textsc{randrange}(0,\, |\hat{r}_{\mathit{rb}}| + 1)$
        \tcp{posición de inserción; $+1$ permite insertar al final}
    }
    $\hat{s}' \leftarrow \textsc{CopiaProfunda}(\hat{s})$
    \tcp{Python pasa listas por referencia; sin copia se corrompería $\hat{s}$}
    $\mathit{bloque} \leftarrow \hat{r}'_{\mathit{ra}}[i\,{:}\,i{+}2]$\;
    $\text{eliminar } \hat{r}'_{\mathit{ra}}[i\,{:}\,i{+}2]$\;
    $\hat{r}'_{\mathit{rb}}[j\,{:}\,j] \leftarrow \mathit{bloque}$
    \tcp{inserción sin sobreescribir: slice vacío $[j:j]$ abre espacio en la posición $j$}
    $m \leftarrow \textsc{MovimientoVecindario}(\texttt{or\_opt\_2},\; \mathit{ruta\_a}=\mathit{ra},\; \mathit{ruta\_b}=\mathit{rb},\; i=i,\; j=j,\; k=2)$\;
    $s' \leftarrow \textsc{RestaurarDepósito}(\hat{s}')$
    \tcp{añade "D" al inicio y fin de cada ruta antes de devolver}
    \Return $(s',\, m)$\;
}
\end{algorithm}

\noindent\textbf{Ejemplo visual} ($i=1$, $j=1$, $r_a \neq r_b$):
\[
    \begin{array}{ll}
        r_a = \bigl[\mathtt{TR1},\;
                    \overbrace{\mathtt{TR2},\;\mathtt{TR3}}^{\mathit{block}},\;
                    \mathtt{TR4}\bigr]
        &
        r_b = [\mathtt{TR5},\;\mathtt{TR6},\;\mathtt{TR7}] \\[8pt]
        \multicolumn{2}{c}{\xrightarrow{\;\text{quitar bloque de }r_a,\;\text{insertar en }r_b[1]\;}} \\[4pt]
        r'_a = [\mathtt{TR1},\;\mathtt{TR4}]
        &
        r'_b = [\mathtt{TR5},\;\mathtt{TR2},\;\mathtt{TR3},\;\mathtt{TR6},\;\mathtt{TR7}]
    \end{array}
\]

\noindent El bloque $[\mathtt{TR2}, \mathtt{TR3}]$ se inserta en su orden
original; no se realiza ninguna inversión.

% ---- or_opt_3 ----
\subsubsection{\texttt{or\_opt\_3}}
\label{sssec:or-opt-3}

\textbf{Movimiento.}
Idéntico a \texttt{or\_opt\_2} pero el bloque tiene exactamente \textbf{3}
tareas consecutivas: $\mathit{block} = r_a[i{:}i{+}3]$.
El bloque se inserta en la posición $j$ de $r_b$ en el \emph{mismo orden}
(sin inversión).  Mover un bloque mayor genera vecinos más disruptivos y
ayuda a la búsqueda a escapar de mínimos locales más profundos.

\textbf{Requisito.} $|r_a| \geq 3$.  El respaldo de ruta única se aplica
igual que en \texttt{or\_opt\_2}: si solo existe una ruta, $r_b$ puede
igualar a $r_a$ y el bloque se reubica dentro de la misma ruta.

\begin{algorithm}[H]
\DontPrintSemicolon
\caption{\texttt{or\_opt\_3} --- Mover un bloque de 3 tareas consecutivas a otra ruta}
\label{alg:or-opt-3}

\KwIn{Solución con depósito $s$;\; generador aleatorio $\mathit{rng}$;\; máximo de intentos $T_{\max}=500$}
\KwOut{Vecino con depósito $s'$;\; registro de movimiento $m$}

\BlankLine
$\hat{s} \leftarrow \textsc{NormalizarQuitarDepósito}(s)$
\tcp{elimina el token "D" de inicio y fin de cada ruta}
$t \leftarrow 0$\;
\Repeat{movimiento válido encontrado}{
    $t \leftarrow t + 1$\;
    \lIf{$t > T_{\max}$}{\textbf{lanzar} \texttt{RuntimeError}}
    $\mathcal{C} \leftarrow \{k : |\hat{r}_k| \geq 3\}$
    \tcp{rutas candidatas como origen: al menos 3 tareas (tamaño del bloque)}
    \lIf{$\mathcal{C} = \emptyset$}{\textbf{continue}}
    $\mathit{ra} \leftarrow \mathit{rng}.\textsc{choice}(\mathcal{C})$\;
    $\mathcal{D} \leftarrow \{x : x \neq \mathit{ra}\}$
    \tcp{rutas destino distintas a la origen; si solo hay una, $\mathcal{D} = \{\mathit{ra}\}$}
    \lIf{$\mathcal{D} = \emptyset$}{$\mathcal{D} \leftarrow \{\mathit{ra}\}$}
    $\mathit{rb} \leftarrow \mathit{rng}.\textsc{choice}(\mathcal{D})$\;
    $n_a \leftarrow |\hat{r}_{\mathit{ra}}|$\;
    $i \leftarrow \mathit{rng}.\textsc{randrange}(0,\, n_a - 3 + 1)$
    \tcp{inicio del bloque; garantiza que caben 3 tareas desde $i$}
    \eIf{$\mathit{rb} = \mathit{ra}$}{
        $\tau \leftarrow n_a - 3$
        \tcp{posiciones disponibles tras extraer el bloque}
        \lIf{$\tau \leq 0$}{\textbf{continue}}
        $j \leftarrow \mathit{rng}.\textsc{randrange}(0,\, \tau + 1)$\;
        \lIf{$j = i$}{\textbf{continue} \tcp*{no genera vecino distinto}}
    }{
        $j \leftarrow \mathit{rng}.\textsc{randrange}(0,\, |\hat{r}_{\mathit{rb}}| + 1)$
        \tcp{posición de inserción; $+1$ permite insertar al final}
    }
    $\hat{s}' \leftarrow \textsc{CopiaProfunda}(\hat{s})$
    \tcp{Python pasa listas por referencia; sin copia se corrompería $\hat{s}$}
    $\mathit{bloque} \leftarrow \hat{r}'_{\mathit{ra}}[i\,{:}\,i{+}3]$\;
    $\text{eliminar } \hat{r}'_{\mathit{ra}}[i\,{:}\,i{+}3]$\;
    $\hat{r}'_{\mathit{rb}}[j\,{:}\,j] \leftarrow \mathit{bloque}$
    \tcp{inserción sin sobreescribir: slice vacío $[j:j]$ abre espacio en la posición $j$}
    $m \leftarrow \textsc{MovimientoVecindario}(\texttt{or\_opt\_3},\; \mathit{ruta\_a}=\mathit{ra},\; \mathit{ruta\_b}=\mathit{rb},\; i=i,\; j=j,\; k=3)$\;
    $s' \leftarrow \textsc{RestaurarDepósito}(\hat{s}')$
    \tcp{añade "D" al inicio y fin de cada ruta antes de devolver}
    \Return $(s',\, m)$\;
}
\end{algorithm}

\noindent\textbf{Ejemplo visual} ($i=0$, $j=1$, $r_a \neq r_b$):
\[
    \begin{array}{ll}
        r_a = \bigl[\overbrace{\mathtt{TR1},\;\mathtt{TR2},\;\mathtt{TR3}}^{\mathit{block}},\;
                    \mathtt{TR4}\bigr]
        &
        r_b = [\mathtt{TR5},\;\mathtt{TR6},\;\mathtt{TR7}] \\[8pt]
        \multicolumn{2}{c}{\xrightarrow{\;\text{quitar bloque de }r_a,\;\text{insertar en }r_b[1]\;}} \\[4pt]
        r'_a = [\mathtt{TR4}]
        &
        r'_b = [\mathtt{TR5},\;\mathtt{TR1},\;\mathtt{TR2},\;\mathtt{TR3},\;\mathtt{TR6},\;\mathtt{TR7}]
    \end{array}
\]

\noindent Como en \texttt{or\_opt\_2}, el bloque $[\mathtt{TR1}, \mathtt{TR2}, \mathtt{TR3}]$
se inserta en su orden original.

% ============================================================
\subsection{Tabla Resumen}
\label{subsec:operator-summary}
% ============================================================

La Tabla~\ref{tab:operator-summary} resume las propiedades clave de los
nueve operadores.

\begin{table}[h]
\centering
\renewcommand{\arraystretch}{1.4}
\begin{tabular}{llccp{5.2cm}}
\hline
\textbf{Operador} & \textbf{Tipo} & \textbf{Tareas mín.} & \textbf{Rutas} & \textbf{Efecto} \\
\hline
\texttt{relocate\_intra}  & intra & 2      & 1 & Mover una tarea a otra posición dentro de la misma ruta \\
\texttt{swap\_intra}      & intra & 2      & 1 & Intercambiar dos tareas dentro de la misma ruta \\
\texttt{2opt\_intra}      & intra & 3      & 1 & Invertir un segmento contiguo dentro de la misma ruta \\
\texttt{relocate\_inter}  & inter & 1      & 2 & Mover una tarea de una ruta a otra (el destino puede estar vacío) \\
\texttt{swap\_inter}      & inter & 1$+$1  & 2 & Intercambiar una tarea entre dos rutas distintas no vacías \\
\texttt{2opt\_star}       & inter & 1$+$1  & 2 & Intercambiar las colas de dos rutas en puntos de corte elegidos \\
\texttt{cross\_exchange}  & inter & 2$+$2  & 2 & Intercambiar segmentos contiguos de $\geq$2 tareas entre dos rutas \\
\texttt{or\_opt\_2}       & inter & 2      & 2$^*$ & Reubicar un bloque de exactamente 2 tareas consecutivas (orden preservado) \\
\texttt{or\_opt\_3}       & inter & 3      & 2$^*$ & Reubicar un bloque de exactamente 3 tareas consecutivas (orden preservado) \\
\hline
\multicolumn{5}{l}{$^*$ Recae en reubicación intra-ruta cuando solo existe una ruta.}
\end{tabular}
\caption{Resumen de los nueve operadores de vecindario.
         ``Tareas mín.'' indica el número mínimo de tareas requerido en la(s)
         ruta(s) de origen; ``Rutas'' indica el número de rutas distintas
         normalmente involucradas.}
\label{tab:operator-summary}
\end{table}
